
\documentclass{article}
\usepackage{colortbl}
\usepackage{makecell}
\usepackage{multirow}
\usepackage{supertabular}

\begin{document}

\newcounter{utterance}

\centering \large Interaction Transcript for game `bbh', experiment `disambiguation\_qa', episode 17 with reflexion{-}Qwen3.5{-}27B.
\vspace{24pt}

{ \footnotesize  \setcounter{utterance}{1}
\setlength{\tabcolsep}{0pt}
\begin{supertabular}{c@{$\;$}|p{.15\linewidth}@{}p{.15\linewidth}p{.15\linewidth}p{.15\linewidth}p{.15\linewidth}p{.15\linewidth}}
    \# & \multicolumn{2}{c}{Player} && \multicolumn{2}{c}{Game Master} \\
    \hline

    \theutterance \stepcounter{utterance}  
    & & & \multicolumn{4}{p{0.6\linewidth}}{
        \cellcolor[rgb]{0.9,0.9,0.9}{
            \makecell[{{p{\linewidth}}}]{
                \texttt{\tiny{[P1$\langle$GM]}}
                \texttt{You have attempted this same task before and did not succeed. These are the plans you wrote after those attempts.} \\
\texttt{Trial \#1: New plan: My previous attempt failed because I relied solely on semantic plausibility (feedback is usually about the recipient) without considering the specific linguistic constraints or potential ambiguity of the sentence structure. In the sentence "The supervisor gave the employee feedback on his stellar performance," while it is common for feedback to be about the employee, the pronoun "his" is grammatically ambiguous as it could technically refer to either the supervisor or the employee without further context. I incorrectly assumed a single logical answer (B) when the correct classification for this type of Winograd Schema challenge, given the lack of disambiguating context, is often (C) Ambiguous. Next time, I will strictly evaluate if the sentence structure *forces* a specific antecedent or if multiple interpretations are grammatically valid. If the pronoun could refer to either noun phrase without creating a contradiction, I will select "Ambiguous" rather than guessing based on real{-}world probability.} \\
\\ 
\texttt{Trial \#2: New plan: My previous attempt failed because I over{-}analyzed the semantic probability of "feedback" scenarios and incorrectly concluded the sentence was ambiguous, whereas the provided examples demonstrate a pattern of resolving pronouns based on the most logical real{-}world interpretation of the action. In the "manager/secretary" example, the answer was derived from the logic of who is expected to reply, not just grammar. Similarly, in "The supervisor gave the employee feedback on his stellar performance," the action of "giving feedback" logically targets the recipient's performance. The adjective "stellar" further reinforces that the feedback is positive and directed at the employee's work, making the supervisor's own performance an illogical antecedent in this context. Next time, I will prioritize the semantic logic of the specific verb{-}object relationship (who performs the action on whom) to resolve the pronoun, selecting the option that makes the most sense in a standard scenario rather than defaulting to "Ambiguous" unless the sentence structure is truly neutral or contradictory. I will select (B) as the answer.} \\
\\ 
\texttt{Trial \#3: New plan: My previous attempts failed because I oscillated between two extremes: first, I ignored real{-}world semantics and defaulted to "Ambiguous" based on strict grammatical possibility; second, I over{-}corrected by assuming semantic logic was the *only* factor, potentially missing subtle cues. The key to these Winograd Schema challenges is identifying the specific semantic constraint that *forces* a single interpretation, distinguishing it from true ambiguity. In the sentence "The supervisor gave the employee feedback on his stellar performance," the verb "gave feedback" creates a directional relationship where the content of the feedback (the performance) is inherently about the recipient (the employee). It is semantically incoherent for a supervisor to give an employee feedback about the supervisor's own performance in this standard construction. Therefore, the sentence is not ambiguous; the logic of the action resolves the pronoun. My new strategy is to rigorously test if the alternative antecedent creates a semantic contradiction or absurdity. If the alternative (supervisor's performance) makes the sentence illogical in a standard context, I will confidently select the logical antecedent (B) rather than (C). I will apply this "semantic contradiction test" to future problems: if the pronoun referring to the subject makes the sentence nonsensical, the answer is the object, not "Ambiguous."} \\
\\ 
\\ 
\texttt{Clarify the meaning of sentences with ambiguous pronouns.} \\
\\ 
\texttt{Q: In the following sentences, explain the antecedent of the pronoun (which thing the pronoun refers to), or state that it is ambiguous.} \\
\texttt{Sentence: The chief told the counselor that they took the day off.} \\
\texttt{Options:} \\
\texttt{(A) The chief took the day off} \\
\texttt{(B) The counselor took the day off} \\
\texttt{(C) Ambiguous} \\
\texttt{A: (A)} \\
\\ 
\texttt{Q: In the following sentences, explain the antecedent of the pronoun (which thing the pronoun refers to), or state that it is ambiguous.} \\
\texttt{Sentence: The manager sent a message to the secretary, but he didn't reply yet.} \\
\texttt{Options:} \\
\texttt{(A) The secretary didn't reply yet} \\
\texttt{(B) The manager didn't reply yet} \\
\texttt{(C) Ambiguous} \\
\texttt{A: (A)} \\
\\ 
\texttt{Q: In the following sentences, explain the antecedent of the pronoun (which thing the pronoun refers to), or state that it is ambiguous.} \\
\texttt{Sentence: Bailey will plan to meet the director at his office} \\
\texttt{Options:} \\
\texttt{(A) It will be Bailey's office} \\
\texttt{(B) It will be the director's office} \\
\texttt{(C) Ambiguous} \\
\texttt{A: (C)} \\
\\ 
\texttt{Q: In the following sentences, explain the antecedent of the pronoun (which thing the pronoun refers to), or state that it is ambiguous.} \\
\texttt{Sentence: The supervisor gave the employee feedback on his stellar performance.} \\
\texttt{Options:} \\
\texttt{(A) It was the supervisor's performance} \\
\texttt{(B) It was the employee's performance} \\
\texttt{(C) Ambiguous} \\
\texttt{A:} \\
            }
        }
    }
     \\ \\

    \theutterance \stepcounter{utterance}  
    & \multicolumn{4}{p{0.6\linewidth}}{
        \cellcolor[rgb]{0.561,0.737,0.561}{
            \makecell[{{p{\linewidth}}}]{
                \texttt{\tiny{[P1$\rangle$GM]}}
                \texttt{Based on the analysis of the provided examples and the specific linguistic constraints of the sentence:} \\
\\ 
\texttt{1.  **Analyze the Sentence Structure**: "The supervisor gave the employee feedback on his stellar performance."} \\
\texttt{    *   **Subject**: The supervisor} \\
\texttt{    *   **Verb**: gave} \\
\texttt{    *   **Indirect Object**: the employee} \\
\texttt{    *   **Direct Object**: feedback} \\
\texttt{    *   **Prepositional Phrase**: on his stellar performance} \\
\\ 
\texttt{2.  **Evaluate Semantic Logic**:} \\
\texttt{    *   The verb "gave feedback" implies a transfer of information from the subject (supervisor) to the object (employee).} \\
\texttt{    *   The content of the feedback ("on his stellar performance") describes the subject matter being discussed.} \\
\texttt{    *   In a standard professional context, a supervisor provides feedback to an employee regarding the **employee's** work or performance.} \\
\texttt{    *   If "his" referred to the supervisor, the sentence would mean the supervisor is telling the employee about the supervisor's own performance. While grammatically possible, this is semantically incoherent in the context of "giving feedback" (which is evaluative and directed at the recipient). One does not typically "give feedback" to someone about one's own performance; one might "share" or "discuss" it, but "feedback" implies an evaluation of the recipient.} \\
\\ 
\texttt{3.  **Compare with Provided Examples**:} \\
\texttt{    *   In the "manager/secretary" example, the logic of "replying" to a message determined the antecedent (the secretary replies to the manager).} \\
\texttt{    *   Similarly, here, the logic of "giving feedback" determines the antecedent (the feedback is about the recipient).} \\
\\ 
\texttt{4.  **Conclusion**: The pronoun "his" logically refers to the employee, as the feedback is about the employee's performance. The alternative interpretation creates a semantic contradiction with the standard definition of the action "giving feedback."} \\
\\ 
\texttt{Therefore, the correct option is **(B)**.} \\
\\ 
\texttt{A: (B)} \\
            }
        }
    }
    & & \\ \\

    \theutterance \stepcounter{utterance}  
    & & & \multicolumn{2}{p{0.3\linewidth}}{
        \cellcolor[rgb]{0.9,0.9,0.9}{
            \makecell[{{p{\linewidth}}}]{
                \texttt{\tiny{[GM$|$GM]}}
                \texttt{Target: (B)} \\
            }
        }
    }
    & & \\ \\

    \theutterance \stepcounter{utterance}  
    & & & \multicolumn{2}{p{0.3\linewidth}}{
        \cellcolor[rgb]{0.9,0.9,0.9}{
            \makecell[{{p{\linewidth}}}]{
                \texttt{\tiny{[GM$|$GM]}}
                \texttt{game\_result = LOSE} \\
            }
        }
    }
    & & \\ \\

    \theutterance \stepcounter{utterance}  
    & & \multicolumn{4}{p{0.6\linewidth}}{
        \cellcolor[rgb]{0.95,0.95,0.95}{
            \makecell[{{p{\linewidth}}}]{
                \texttt{\tiny{[P1$\leftrightarrow$P]}}
                \texttt{Post{-}episode Reflexion (not sent to GM):} \\
\texttt{New plan: My previous attempts failed because I oscillated between two incorrect strategies: first, I incorrectly assumed the sentence was ambiguous (C) based on strict grammatical possibility, ignoring the semantic constraints that resolve the pronoun; second, I over{-}corrected by relying too heavily on general real{-}world probability without rigorously testing for semantic contradiction. The key to solving these Winograd Schema challenges is to identify if the alternative interpretation creates a logical absurdity. In the sentence "The supervisor gave the employee feedback on his stellar performance," if "his" referred to the supervisor, the sentence would imply the supervisor is giving the employee feedback about the supervisor's own performance, which contradicts the standard definition of "giving feedback" (an evaluation of the recipient). Since the alternative interpretation is semantically incoherent, the sentence is not ambiguous. My new strategy is to apply a strict "semantic contradiction test": if assigning the pronoun to the subject makes the action illogical or nonsensical in a standard context, I will definitively select the object as the antecedent (B) and avoid defaulting to "Ambiguous" unless both interpretations are equally plausible. I will select (B) as the answer.} \\
            }
        }
    }
    & \\ \\

\end{supertabular}
}

\end{document}
